\chapter{Chatting with a Stranger on Your Trip in Japan}

% ================= PERCAKAPAN 1-4 =================
\begin{convobox}{1}
    \noindent
    \jpkana{こんにちは。}{Konnichiwa.} 
    \jpword{きょう}{今日}{Kyou} \jpkana{は}{wa} 
    \jpword{りょこう}{ご旅行}{go-ryokou} \jpkana{ですか？}{desu ka?}

    \vspace{1.5em}
    \noindent{\Large \color{articolor} \textbf{Arti:} Halo. Apakah hari ini kamu sedang jalan-jalan?}
\end{convobox}

\begin{convobox}{2}
    \noindent
    \jpkana{はい、}{Hai,} 
    \jpword{りょこう}{旅行}{ryokou} \jpkana{です。}{desu.}

    \vspace{1.5em}
    \noindent{\Large \color{articolor} \textbf{Arti:} Ya, (sedang) jalan-jalan.}
\end{convobox}

\begin{convobox}{3}
    \noindent
    \jpkana{どちら}{Dochira} \jpkana{の}{no} 
    \jpword{くに}{国}{kuni} \jpkana{から}{kara} 
    \jpkana{いらっしゃった}{irasshatta} \jpkana{ん}{n} \jpkana{ですか？}{desu ka?}

    \vspace{1.5em}
    \noindent{\Large \color{articolor} \textbf{Arti:} Dari negara mana kamu berasal?}
\end{convobox}

\begin{convobox}{4}
    \noindent
    \jpkana{インドネシア}{Indoneshia} \jpkana{から}{kara} 
    \jpword{き}{来}{ki}\jpkana{ました。}{mashita.}

    \vspace{1.5em}
    \noindent{\Large \color{articolor} \textbf{Arti:} Saya datang dari Indonesia.}
\end{convobox}

% --- LESSON BUFFER 1 ---
\begin{lessonbox}{Catatan Belajar: Asal Negara!}
    Halo teman-teman! Hebat banget kan Andi berani ngobrol sama orang Jepang asli? Di kalimat 3 dan 4, kita belajar cara kasih tahu dari mana kita berasal.\\
    
    Kalau ada yang tanya pakai kata \textbf{"\dots kara kimashita"} (datang dari\dots), kamu tinggal sebut negaramu lalu ditambah kata itu. Contohnya: \textit{Indoneshia kara kimashita!} Gampang banget, kan? Yuk lanjut lihat mereka mau ke mana!
\end{lessonbox}

% ================= PERCAKAPAN 5-9 =================
\begin{convobox}{5}
    \noindent
    \jpword{にほん}{日本}{Nihon} \jpkana{は}{wa} 
    \jpword{はじ}{初}{haji}\jpkana{めて}{mete} \jpkana{ですか？}{desu ka?}

    \vspace{1.5em}
    \noindent{\Large \color{articolor} \textbf{Arti:} Apakah ini pertama kalinya ke Jepang?}
\end{convobox}

\begin{convobox}{6}
    \noindent
    \jpkana{はい、}{Hai,} 
    \jpword{はじ}{初}{haji}\jpkana{めて}{mete} \jpkana{です。}{desu.}

    \vspace{1.5em}
    \noindent{\Large \color{articolor} \textbf{Arti:} Ya, pertama kali.}
\end{convobox}

\begin{convobox}{7}
    \noindent
    \jpword{きょう}{今日}{Kyou} \jpkana{は}{wa} 
    \jpkana{どちら}{dochira} \jpkana{まで}{made} 
    \jpword{い}{行}{i}\jpkana{かれる}{kareru} \jpkana{ん}{n} \jpkana{ですか？}{desu ka?}

    \vspace{1.5em}
    \noindent{\Large \color{articolor} \textbf{Arti:} Hari ini mau pergi sampai ke mana?}
\end{convobox}

\begin{convobox}{8}
    \noindent
    \jpword{きょうと}{京都}{Kyouto} \jpkana{まで}{made} 
    \jpword{い}{行}{i}\jpkana{きます。}{kimasu.}

    \vspace{1.5em}
    \noindent{\Large \color{articolor} \textbf{Arti:} Pergi sampai ke Kyoto.}
\end{convobox}

\begin{convobox}{9}
    \noindent
    \jpkana{そうですか。}{Sou desu ka.} 
    \jpword{わたし}{私}{Watashi} \jpkana{は}{wa} 
    \jpword{なごや}{名古屋}{Nagoya} \jpkana{まで}{made} 
    \jpword{い}{行}{i}\jpkana{きます。}{kimasu.}

    \vspace{1.5em}
    \noindent{\Large \color{articolor} \textbf{Arti:} Oh begitu. Saya pergi sampai ke Nagoya.}
\end{convobox}

% --- LESSON BUFFER 2 ---
\begin{lessonbox}{Catatan Belajar: Mau ke mana kita?}
    Perhatikan kata \textbf{まで (made)} di atas! Kata ini artinya "sampai". Kalau kamu mau bilang pergi sampai ke suatu tempat, sebut nama kotanya ditambah \textit{made ikimasu}. \\
    
    Orang Jepang suka nanya \textit{"dochira made?"} (mau sampai ke arah mana?). Nah, sekarang siapkan tiketmu, karena Andi mau ke Kyoto dan teman barunya mau ke Nagoya!
\end{lessonbox}

% ================= PERCAKAPAN 10-17 =================
\begin{convobox}{10}
    \noindent
    \jpword{とうちゃく}{到着}{Touchaku} \jpkana{したら、}{shitara,} 
    \jpkana{どこ}{doko} \jpkana{に}{ni} 
    \jpword{い}{行}{i}\jpkana{かれる}{kareru} 
    \jpword{よてい}{予定}{yotei} \jpkana{ですか？}{desu ka?}

    \vspace{1.5em}
    \noindent{\Large \color{articolor} \textbf{Arti:} Kalau sudah tiba, rencana mau pergi ke mana?}
\end{convobox}

\begin{convobox}{11}
    \noindent
    \jpword{さくら}{桜}{Sakura} \jpkana{が}{ga} 
    \jpword{ゆうめい}{有名}{yuumei} \jpkana{な}{na} 
    \jpword{てら}{お寺}{o-tera} \jpkana{に}{ni} 
    \jpword{い}{行}{i}\jpkana{く}{ku} \jpword{よてい}{予定}{yotei} \jpkana{です。}{desu.}

    \vspace{1.5em}
    \noindent{\Large \color{articolor} \textbf{Arti:} Rencananya pergi ke kuil yang terkenal dengan bunga sakuranya.}
\end{convobox}

\begin{convobox}{12}
    \noindent
    \jpword{たの}{楽}{Tano}\jpkana{しそう}{shisou} \jpkana{ですね。}{desu ne.}

    \vspace{1.5em}
    \noindent{\Large \color{articolor} \textbf{Arti:} Sepertinya menyenangkan ya.}
\end{convobox}

\begin{convobox}{13}
    \noindent
    \jpkana{それにしても、}{Sore ni shitemo,} 
    \jpword{にほんご}{日本語}{Nihongo} \jpkana{が}{ga} 
    \jpkana{とても}{totemo} \jpword{じょうず}{上手}{jouzu} \jpkana{ですね。}{desu ne.}

    \vspace{1.5em}
    \noindent{\Large \color{articolor} \textbf{Arti:} Ngomong-ngomong, bahasa Jepangnya sangat pintar ya.}
\end{convobox}

\begin{convobox}{14}
    \noindent
    \jpkana{いつ}{Itsu} \jpkana{から}{kara} 
    \jpword{にほんご}{日本語}{Nihongo} \jpkana{を}{o} 
    \jpword{べんきょう}{勉強}{benkyou} \jpkana{している}{shite iru} \jpkana{ん}{n} \jpkana{ですか？}{desu ka?}

    \vspace{1.5em}
    \noindent{\Large \color{articolor} \textbf{Arti:} Sejak kapan belajar bahasa Jepang?}
\end{convobox}

\begin{convobox}{15}
    \noindent
    \jpword{いちねんまえ}{1年前}{Ichi-nen mae} \jpkana{から}{kara} 
    \jpword{べんきょう}{勉強}{benkyou} \jpkana{しています。}{shite imasu.}

    \vspace{1.5em}
    \noindent{\Large \color{articolor} \textbf{Arti:} Belajar sejak 1 tahun yang lalu.}
\end{convobox}

\begin{convobox}{16}
    \noindent
    \jpkana{なんで}{Nande} \jpword{にほんご}{日本語}{Nihongo} \jpkana{の}{no} 
    \jpword{べんきょう}{勉強}{benkyou} \jpkana{を}{o} 
    \jpword{はじ}{始}{haji}\jpkana{めた}{meta} \jpkana{ん}{n} \jpkana{ですか？}{desu ka?}

    \vspace{1.5em}
    \noindent{\Large \color{articolor} \textbf{Arti:} Kenapa mulai belajar bahasa Jepang?}
\end{convobox}

\begin{convobox}{17}
    \noindent
    \jpword{にほん}{日本}{Nihon} \jpkana{の}{no} 
    \jpword{ぶんか}{文化}{bunka} \jpkana{と}{to} 
    \jpword{れきし}{歴史}{rekishi} \jpkana{に}{ni} 
    \jpword{きょうみ}{興味}{kyoumi} \jpkana{が}{ga} \jpkana{ある}{aru} \jpkana{から}{kara} \jpkana{です。}{desu.}

    \vspace{1.5em}
    \noindent{\Large \color{articolor} \textbf{Arti:} Karena saya tertarik dengan budaya dan sejarah Jepang.}
\end{convobox}

% --- LESSON BUFFER 3 ---
\begin{lessonbox}{Catatan Belajar: Pujian dan Hobi}
    Wah, bahasa Jepang Andi dipuji \textbf{上手 (jouzu)} yang artinya pintar/jago! Kalau kamu dipuji seperti ini, kamu bisa senyum bahagia lho. \\
    
    Di sini Andi juga ngomongin alasannya belajar bahasa Jepang karena tertarik dengan \textit{bunka} (budaya) dan \textit{rekishi} (sejarah). Orang Jepang sangat senang loh kalau kita tertarik dengan negara mereka!
\end{lessonbox}

% ================= PERCAKAPAN 18-27 =================
\begin{convobox}{18}
    \noindent
    \jpword{にほん}{日本}{Nihon} \jpkana{の}{no} 
    \jpword{た}{食}{ta}\jpkana{べ}{be}\jpword{もの}{物}{mono} \jpkana{は}{wa} 
    \jpword{す}{好}{su}\jpkana{き}{ki} \jpkana{ですか？}{desu ka?}

    \vspace{1.5em}
    \noindent{\Large \color{articolor} \textbf{Arti:} Apakah kamu suka makanan Jepang?}
\end{convobox}

\begin{convobox}{19}
    \noindent
    \jpkana{はい、}{Hai,} \jpword{だいす}{大好}{daisu}\jpkana{き}{ki} \jpkana{です。}{desu.}

    \vspace{1.5em}
    \noindent{\Large \color{articolor} \textbf{Arti:} Ya, sangat suka.}
\end{convobox}

\begin{convobox}{20}
    \noindent
    \jpword{いま}{今}{Ima} \jpkana{まで}{made} 
    \jpword{にほん}{日本}{Nihon} \jpkana{で}{de} 
    \jpword{た}{食}{ta}\jpkana{べた}{beta} \jpword{もの}{物}{mono} \jpkana{の}{no} 
    \jpword{なか}{中}{naka} \jpkana{で、}{de,}

    \vspace{1.5em}
    \noindent{\Large \color{articolor} \textbf{Arti:} Di antara makanan yang pernah dimakan di Jepang selama ini,}
\end{convobox}

\begin{convobox}{21}
    \noindent
    \jpword{なに}{何}{Nani} \jpkana{が}{ga} 
    \jpword{いちばん}{一番}{ichiban} \jpkana{おいしかった}{oishikatta} \jpkana{ですか？}{desu ka?}

    \vspace{1.5em}
    \noindent{\Large \color{articolor} \textbf{Arti:} Apa yang paling lezat?}
\end{convobox}

\begin{convobox}{22}
    \noindent
    \jpword{とうきょう}{東京}{Toukyou} \jpkana{で}{de} 
    \jpword{た}{食}{ta}\jpkana{べた}{beta} \jpword{すし}{寿司}{sushi} \jpkana{が}{ga} 
    \jpword{いちばん}{一番}{ichiban} \jpkana{おいしかった}{oishikatta} \jpkana{です。}{desu.}

    \vspace{1.5em}
    \noindent{\Large \color{articolor} \textbf{Arti:} Sushi yang saya makan di Tokyo adalah yang paling lezat.}
\end{convobox}

\begin{convobox}{23}
    \noindent
    \jpkana{（まもなく、}{Mamonaku,} 
    \jpword{なごや}{名古屋}{Nagoya}\jpkana{、}{,} \jpword{なごや}{名古屋}{Nagoya} \jpkana{です。）}{desu.}

    \vspace{1.5em}
    \noindent{\Large \color{articolor} \textbf{Arti:} (Sebentar lagi, tiba di Nagoya, Nagoya.)}
\end{convobox}

\begin{convobox}{24}
    \noindent
    \jpkana{あっ、}{A,,} \jpkana{もう}{mou} 
    \jpword{つ}{着}{tsu}\jpkana{きました。}{kimashita.}

    \vspace{1.5em}
    \noindent{\Large \color{articolor} \textbf{Arti:} Ah, sudah sampai.}
\end{convobox}

\begin{convobox}{25}
    \noindent
    \jpword{しんかんせん}{新幹線}{Shinkansen} \jpkana{は}{wa} 
    \jpword{はや}{速}{haya}\jpkana{い}{i} \jpkana{ですね。}{desu ne.}

    \vspace{1.5em}
    \noindent{\Large \color{articolor} \textbf{Arti:} Shinkansen itu cepat ya.}
\end{convobox}

\begin{convobox}{26}
    \noindent
    \jpword{にほん}{日本}{Nihon} \jpword{りょこう}{旅行}{ryokou}\jpkana{、}{,} 
    \jpword{たの}{楽}{tano}\jpkana{しんで}{shinde} \jpkana{ください。}{kudasai.}

    \vspace{1.5em}
    \noindent{\Large \color{articolor} \textbf{Arti:} Tolong nikmati perjalananmu di Jepang ya.}
\end{convobox}

\begin{convobox}{27}
    \noindent
    \jpkana{それでは、}{Sore dewa,} 
    \jpword{しつれい}{失礼}{shitsurei} \jpkana{します。}{shimasu.}

    \vspace{1.5em}
    \noindent{\Large \color{articolor} \textbf{Arti:} Kalau begitu, saya permisi dulu.}
\end{convobox}

% --- SUMMARY BOX ---
\vspace{2em}
\begin{tcolorbox}[
    colback=blue!5!white,
    colframe=blue!80!black,
    boxrule=3pt,
    arc=5mm,
    title={\huge\bfseries \sffamily 🌟 Rangkuman Bab 1 🌟},
    fonttitle=\color{white},
    attach boxed title to top center={yshift=-4mm},
    boxed title style={colback=blue!80!black, rounded corners},
    left=5mm, right=5mm, top=7mm, bottom=5mm
]
\Large \textbf{Kerja bagus teman-teman!} Di bab ini kita sudah belajar banyak hal keren saat ngobrol di dalam kereta Jepang (Shinkansen):
\begin{itemize}
    \item \textbf{Menjawab asal negara:} "\textit{\dots kara kimashita}" (Datang dari\dots)
    \item \textbf{Bilang ke mana kita pergi:} "\textit{\dots made ikimasu}" (Pergi sampai ke\dots)
    \item \textbf{Bicara soal durasi:} "\textit{Ichi-nen mae kara\dots}" (Sejak 1 tahun lalu\dots)
    \item \textbf{Membicarakan kesukaan (makanan):} "\textit{Daisuki desu!}" (Sangat suka!)
    \item \textbf{Berpamitan sopan:} "\textit{Shitsurei shimasu}" (Permisi / Saya pamit)
\end{itemize}
Jangan lupa diulang-ulang ya bacanya biar makin lancar seperti Andi! Sampai jumpa di perjalanan berikutnya!
\end{tcolorbox}

\newpage
